\documentclass[UTF8]{ctexart}
\usepackage{geometry}
\usepackage{booktabs}
\usepackage{graphicx}
\usepackage{amsmath}
\usepackage{multicol}
\usepackage{listings}
\usepackage{xcolor}
\usepackage{tcolorbox}
\usepackage{subcaption}
\usepackage{abstract}
\usepackage{tabularray}
\usepackage{enumitem}

\renewcommand{\abstractname}{\Large 摘要}
\renewcommand{\abstracttextfont}{\normalfont\fontsize{11pt}{12pt}\selectfont}
\setlength{\absleftindent}{0pt}
\setlength{\absrightindent}{0pt}
\bibliographystyle{plain}
\lstset{
	basicstyle=\ttfamily\small,
	breaklines=true,
	backgroundcolor=\color{gray!10},
	frame=single
}

\title{基于混合整数规划的电网电力智能调控模型}
\author{}
\date{}

\geometry{a4paper,margin=2.5cm}
\begin{document}
	
	\maketitle
	\pagestyle{empty}
	\begin{abstract}
		本文针对微网与外部电网的电力调控问题，按照“确定性优化—不确定性预测—滚动调整—动态电价”的思路，
		建立了由\textbf{混合整数线性规划}、\textbf{时间序列滚动预测}和\textbf{场景随机优化}组成的购电与储能调度模型，
		为微网在不同信息条件和电价机制下制定经济、可靠的购电策略。\par
		\indent \textbf{针对问题一}，将一天划分为144个10分钟时段，以购电量、充放电量、储电量和充放电状态为决策变量，
		以全天购电费用最小为目标，建立\textbf{混合整数线性规划（MILP）}模型。模型同时约束逐时供需平衡、储电量递推、
		1200--10800 kWh储电量范围、5000 kW充放电功率上限、充放电互斥以及日初和日末储电量均为6000 kWh。
		使用HiGHS求解器求得最优计划：全天购电量约为59482.699 kWh，购电费约为35126.949元，
		全天充电量和放电量分别约为20740.666 kWh和16799.940 kWh，储能损耗约为3940.727 kWh。\par
		\indent \textbf{针对问题二}，在负载和光伏均存在不确定性的条件下，建立历史数据驱动的日前购电模型。
		负载预测采用31天滚动岭回归，融合近期滞后值、均值、星期周期和年内周期特征；光伏预测采用扩展窗口岭回归，
		引入太阳赤纬等物理特征，并将其与最近7日同时刻均值按权重融合。基于上述模型生成
		2025年2月1日至12月31日共334天、每天144个时段的预测数据。购电端采用场景随机优化，
		所有场景共用0:00确定的一条计划购电曲线，以计划购电费用、五倍电价紧急购电费用的期望值和日末储能价值之和最小为目标；
		实际运行时储能根据已观测供需逐时段调整，并通过$S_{d+1,0}=S_{d,144}$保证跨日储电量连续。
		单时段无储能的简化分析表明，最优计划购电量对应净负载的80\%分位数，体现了高额紧急购电价格下的风险补偿思想。\par
		\indent \textbf{针对问题三}，利用每天0:00、6:00、12:00和18:00发布的未来24小时光伏预报建立多阶段滚动优化模型。
		以0:00计划购电量为基准，在后续预报时刻根据最新信息调整购电策略；计划量高于调整量的部分按实时电价的50\%支付违约费用，
		调整量高于计划量的部分按实时电价的1.5倍计费。模型在满足紧急购电和储能运行约束的同时，
		优化各预报时刻的调整量，从而降低预测更新带来的费用风险。\par
		\indent \textbf{针对问题四}，将实时波动电价纳入不确定性框架，建立电价与负载、光伏联合滚动预测下的随机优化模型。
		在问题二和问题三的购电机制上分别引入动态电价场景，重新计算计划购电、调整购电、紧急购电及储能充放电策略，
		比较不同预测更新频率和风险参数下的经济性，为实时电价环境中的微网调度提供决策依据。\par
		\noindent \textbf{关键词}：微网调度；混合整数线性规划；滚动岭回归；场景随机优化；储能；动态电价
	\end{abstract}
	
	\clearpage
	
	\pagestyle{plain}
	\section{问题重述}
	微网是一种将光伏发电、小区负载和储能设备等分布式资源进行本地协调的小型发用电系统。
	在发展非孤岛式微网时，系统既需要保证向小区负载持续供电，又需要利用储能设备在外网电价较低或
	光伏发电有剩余时为电池充电，并在外网电价较高时释放电能，从而降低微网的购电费用。因此，
	如何根据电价、小区负载、光伏发电功率和储能设备当前储电量，合理安排各时段的计划购电量与
	储能充放电量，是微网与外部电网协同调控的关键。\par
	题目中的储能设备最大容量为12000 kWh，最大充放电功率为5000 kW，
	2025年1月1日0:00的初始储电量为6000 kWh。为保证设备安全运行，
	后续过程中的储电量必须始终保持在1200--10800 kWh之间，充放电效率均为90\%。\par
	根据附件提供的数据，本题需要依次解决以下四个问题：
	\begin{enumerate}
		\item [(1)] 若每天的电价和小区负载相同，并已知光伏发电功率一天的预测数据，
		要求微网提供的电能不低于小区负载，且储能设备在0:00和24:00的储电量相同。
		需要在每天0:00制定当天的计划购电策略，并给出指定 10 分钟时段的购电量、
		全天购电量和购电费，以及六个四小时区间的充放电量和起止时刻储电量。
		
		\item [(2)] 若每天的电价相同，但小区负载和光伏发电功率随时间变化，
		微网提供的电能仍不得低于小区负载。当实际供电不足时，需要以交易时刻电价的5倍进行紧急购电；
		除紧急购电外，其他计划购电费用仍按计划购电量计算。需要利用附件1中的电价和附件2中的
		实际负载、光伏数据，在每天0:00制定当天的计划购电策略，并给出指定日期的购电量、
		充放电量、储电量及紧急购电结果。
		
		\item [(3)] 在每天0:00、6:00、12:00和18:00可以获得未来24小时整点的光伏发电功率预报。
		系统既可以根据0:00的预报制定全天计划购电策略，也可以根据后续时刻的新预报调整购电策略。
		若最终计划购电量高于调整后的购电量，超出部分按交易时刻电价的50\%支付违约费用；
		若调整后的购电量高于计划购电量，超出部分按交易时刻电价的1.5倍计费。
		需要建立包含计划购电费用、紧急购电费用和调整费用的优化模型，并分析是否需要引入其他时刻的预报。
		
		\item [(4)] 考虑外网电价也随实际交易时刻实时波动。需要利用附件2中的负载和光伏数据、
		附件3中的光伏预报数据以及附件4中的动态电价数据，在波动电价条件下重新建立微网购电模型，
		并分别重新计算问题2和问题3的结果。
	\end{enumerate}
	\hspace{2em}附件1给出了某一天的电价、小区负载和光伏发电预测功率；附件2给出了2025年全年每天
	10分钟粒度的实际小区负载和光伏发电功率；附件3给出了每天0:00、6:00、12:00和18:00发布的
	未来24小时整点光伏预报；附件4给出了全年实时波动的电价数据；附件5给出了各问结果文件模板。
	本文需要综合上述数据和储能运行约束，建立微网购电与储能调度模型，使系统在满足负载需求的前提下
	尽可能降低购电费用。
	\section{问题分析}
	本文要求建立数学模型研究微网与外部电网之间的电力调控策略。核心任务是在储能设备运行约束下，
	根据电价、小区负载、光伏发电功率和当前储电量，合理确定计划购电量及储能充放电量，
	使微网既能满足小区负载需求，又能尽可能降低购电费用。四个问题由确定性调度逐步过渡到
	不确定性预测、滚动调整和动态电价，因此本文总体按照“确定性优化—时间序列预测—场景优化—
	多阶段调整”的思路展开。\par
	\subsection{问题一的分析}
	问题一是电价、小区负载和光伏预测均已知的确定性日前调度问题。其难点在于储能设备将不同
	时段相互耦合：当前充放电不仅影响本时段购电量，还会改变后续时段的剩余储能空间；同时，
	同一时段不能既充电又放电，因此模型含有离散决策。为得到全天费用最小且满足物理约束的
	计划，我们将问题一的求解过程划分为以下几个步骤：
	\begin{enumerate}[label=(\arabic*),leftmargin=2.8em]
		\item 将一天划分为144个10分钟时段，把附件1中的平均功率换算为各时段的负载电量和光伏电量，
		并统一采用电池外部电量口径描述充放电量。
		\item 设置购电量、充电量、放电量、时段末储电量和充放电状态五类决策变量，
		以全天购电费用最小为目标，建立混合整数线性规划模型。
		\item 依次引入逐时供需平衡、储电量递推、储电量上下限、最大充放电功率、充放电互斥以及
		日初日末储电量相等等约束，保证调度计划满足负载需求和储能设备的安全运行条件。
		\item 利用HiGHS求解器求得MILP最优解，并通过独立计算供需残差、储电量递推残差和功率约束残差，
		检验所得计划的可执行性，最后生成表1、表2和\texttt{result1.xlsx}。
	\end{enumerate}
	\subsection{问题二的分析}
	问题二中每天的电价规律保持不变，但小区负载和光伏发电功率逐日、逐时段变化。系统在每天
	0:00制定计划时，不能提前获得当天未来的实际负载和光伏数据，只能利用历史信息进行预测。
	如果按照预测结果制定的计划购电量与实际光伏不能满足实际负载，则先调用储能放电；仍有缺口时，
	再按交易时刻正常电价的5倍进行紧急购电。因此，第二问的本质是一个包含预测不确定性、
	风险修正和储能调度的序贯决策问题，其主要决策链为
	\begin{equation}
		\begin{aligned}
			\text{历史数据}
			&\rightarrow \text{负载与光伏滚动预测}
			\rightarrow \text{预测误差风险修正}\\
			&\rightarrow \text{日前购电与储能优化}
			\rightarrow \text{实际运行与紧急购电}.
		\end{aligned}
	\end{equation}
	为了兼顾提前购电成本和缺电风险，本文将问题二分为以下几个步骤：
	\begin{enumerate}[label=(\arabic*),leftmargin=2.8em]
		\item 对附件2中的负载和光伏数据按日期、时段进行对齐，将每天144个10分钟功率点组成
		一条时间序列。利用2025年1月数据进行模型参数选择和储能热启动，将2025年2月1日至
		12月31日作为正式预测和费用统计区间。
		\item 针对小区负载建立31天滚动岭回归预测模型。模型同时使用前1、2、3、7日的滞后值、
		前3日和7日的同时刻均值、最近两个相同星期几的均值，以及星期周期、年内周期和日内周期特征。
		利用1月验证集选择岭参数，使模型能够适应近期负载水平的变化。
		\item 针对光伏发电建立扩展窗口岭回归预测模型。模型引入前1、2、3、7、14日的滞后值，
		前3日、7日和14日的同时刻均值，以及由太阳赤纬描述的年内季节特征和日内周期特征；
		随后将岭回归结果与最近7日同时刻均值进行加权融合，并对结果进行非负截断。1月验证得到的
		当前参数为岭参数1、融合权重0.4。
		\item 根据负载和光伏预测计算净负荷预测值，并利用最近60天同一时段的滚动预测误差确定
		80\%分位数风险裕量。将该裕量加入预测净负荷，得到计划净负荷。由于实际缺口需要按5倍
		电价紧急购电，80\%分位数修正能够在计划购电成本和缺电风险之间建立量化平衡，避免仅按
		预测均值安排购电。
		\item 每天0:00根据计划净负荷求解一个储能调度MILP，以全天计划购电费用最小为目标。
		模型包括计划购电量、充放电量、计划弃电量、储电量和充放电状态变量，并考虑计划净负荷平衡、
		储电量递推、储电量上下限、充放电功率上限、充放电互斥以及日初日末储电量相等约束。
		\item 将MILP求得的计划购电曲线在当天0:00锁定。实际运行阶段使用附件2中的真实负载和
		光伏逐时段运行：当计划电量和光伏有富余时优先给电池充电，未消纳部分记为计划弃电；
		当出现缺口时优先调用电池放电，电池放电后仍不能补足的差值再进行紧急购电。
		\item 计算每天的实际充电量、放电量、计划购电费用和紧急购电费用，并将当天24:00的
		实际储电量传递给次日0:00，保证储能状态跨日连续。按照上述流程逐日滚动求解至2025年
		12月31日，最终汇总计划购电量、实际充放电量、紧急购电时段和总费用，并生成
		\texttt{result2.xlsx}。
	\end{enumerate}
	\subsection{问题三的分析}
	问题三在第二问的基础上增加了日内光伏预报更新和购电调整机制。每天有四个预报发布与决策时刻
	0:00、6:00、12:00和18:00。0:00根据首版预报制定全天计划购电量，后续三个时刻根据新预报、
	已经发生的实际数据以及当前储电量，只调整尚未结束时段的购电量；实际运行时若最终调整购电量和
	实际光伏仍不能满足负载，则由储能放电，剩余缺口再按5倍电价紧急购电。因此，第三问由单一的
	“计划购电+紧急购电”扩展为“计划购电+多次调整购电+紧急购电”的序贯决策问题。为刻画信息更新
	和费用结算过程，我们将问题三分为以下几个步骤：
	\begin{enumerate}[label=(\arabic*),leftmargin=2.8em]
		\item 明确三类数据的使用方式：负载预测采用时间序列模型；光伏预测只采用附件3在相应决策
		时刻发布的未来24小时整点预报；实际费用和储能运行则使用附件2中的实际负载与实际光伏。
		\item 附件3给出一小时粒度的光伏预报，而购电和储能决策需要细化到10分钟。对相邻整点预报
		进行线性插值，将小时预报展开为144个10分钟时段，并进行两端水平外推，得到各决策时刻
		可使用的完整日内光伏曲线。
		\item 在0:00根据负载预测、首版光伏预报和预测净负荷求解全天计划MILP，确定计划购电量、
		储能充放电量、富余量和计划储电量；此时确定的计划购电曲线全天锁定，已经支付的计划电量
		即使未全部使用也不退款。
		\item 在6:00、12:00和18:00根据最新光伏预报、更新后的负载预测和当前实际储电量重新优化
		未来时段。已经过去或已经执行的时段不可修改，只调整当前时刻之后的购电量，从而形成不回看、
		不使用未来信息的滚动优化过程。
		\item 将每个时段最终调整购电量相对0:00计划量的差额分解为调高量和调低量：调高部分按
		1.5倍实时电价收费，调低部分按0.5倍实时电价收取违约费用，避免对同一时段重复计费。
		\item 在实际运行阶段使用最终调整购电量、实际光伏和实际负载，优先利用储能消纳富余电量，
		缺电时优先放电，放电后仍不足的差值再按5倍实时电价进行紧急购电，并保证储电量跨日连续。
		\item 比较0:00、6:00、12:00和18:00四时刻预报方案与增设其他预报时刻后的费用和缺电风险，
		判断新增预报是否能带来足够的经济改善，从而回答题目中“是否需要引入其他时刻预报”的问题。
	\end{enumerate}
	\subsection{问题四的分析}
	问题四进一步考虑外网电价实时波动，即无论是日前计划还是后续调整，都无法预先准确知道
	未来所有交易时刻的电价。因此，负载、光伏和电价同时构成不确定源，问题由单一的“购电量优化”
	转变为“购电量、调整量和储能动作”的联合决策。为了在动态电价下重新计算问题2和问题3，
	我们将问题四分为以下几个步骤：
	\begin{enumerate}[label=(\arabic*),leftmargin=2.8em]
		\item 对附件4中的实时电价进行整理，建立电价时间序列预测模型，并将价格预测误差与
		负载、光伏预测误差统一表示为随机场景。
		\item 在问题2的框架下，将所有场景中共用的日前计划购电量作为第一阶段决策，
		将实际运行时的储能充放电和紧急购电作为第二阶段决策，以总购电费用的条件期望最小为目标。
		\item 在问题3的框架下，将电价作为各预报时刻更新信息的一部分，重新确定计划购电量、
		调整购电量和费用，建立动态电价条件下的多阶段随机滚动优化模型。
		\item 分别计算动态电价下问题2和问题3的结果，比较不同电价波动幅度和风险偏好对购电费用、
		紧急购电量和储能调度的影响，并生成\texttt{result4-2.xlsx}和\texttt{result4-3.xlsx}。
	\end{enumerate}
	\section{模型假设}
	\begin{enumerate}
		\item 假设附件中的电价、小区负载和光伏发电功率数据真实可靠。各时段的功率数据视为该
		10分钟区间内的平均功率，区间电量为功率与时段长度$\Delta t=1/6$小时的乘积。
		\item 题目所给时间标签表示对应时段的结束时刻，例如附件中的0:10表示0:00--0:10时段，
		\texttt{0:00+1}表示当天24:00。
		\item 假设外网能够满足微网的计划购电和紧急购电需求，不考虑外网故障、限电或供电中断。
		小区负载必须全部得到满足。
		\item 储能设备的储电量采用电池内部电量口径，充入和放出的电量采用电池外部计量口径。
		充电效率和放电效率均为90\%，模型不考虑自放电以及电池老化、循环寿命和折旧费用。
		\item 储能设备的最大充放电功率均为5000 kW，储电量始终保持在1200--10800 kWh之间，
		2025年1月1日0:00的初始储电量为6000 kWh。
		\item 同一10分钟时段内储能设备最多只能充电、放电或闲置，不能同时进行充电和放电；
		每个时段内的平均充放电功率不超过设备功率上限。
		\item 问题一中光伏预测电量全部用于供应小区负载或为储能设备充电，不向外网售电，
		也不额外设置弃光变量；问题二至问题四中未被利用的富余电量不能获得退款或售电收益。
		\item 计划购电策略必须在每天0:00根据当时已经掌握的信息一次性确定，之后不能使用当天
		尚未发生的负载、光伏或电价数据。实际运行时储能设备可以根据已经观测到的供需情况调整
		充放电量；问题三和问题四可在题目规定的预报时刻对后续购电量进行调整。
		\item 紧急购电只在储能调节后仍然存在供电缺口时发生，其电价按交易时刻正常电价的5倍计算，
		且不考虑额外手续费、容量费用或违约以外的其他惩罚费用。
		\item 问题一中储能设备在0:00和24:00的储电量相同；问题二至问题四中储能状态跨日连续，
		即次日0:00的储电量等于当天24:00的储电量。
		\item 时间序列预测模型所使用的历史数据能够反映负载和光伏的基本变化规律，预测误差
		可由历史残差场景近似；模型不考虑通信延迟、设备故障、人工操作误差以及题目未给出的
		其他随机因素。
	\end{enumerate}
	\clearpage
	
	\section{符号说明}
	\begin{table}[h]
		\centering
		\caption{重要符号表}
		\renewcommand{\arraystretch}{1.2}
		\begin{tabular}{clc}
			\toprule
			符号 & 定义说明 & 单位 \\
			\midrule
			$t$ & 时段索引，$t=1,2,\ldots,144$ & -- \\
			$\Delta t$ & 每个时段的长度 & h \\
			$p_t$ & 第$t$时段的购电价格 & 元/kWh \\
			$L_t$ & 第$t$时段小区负载平均功率 & kW \\
			$\widehat{P}_t$ & 第$t$时段光伏发电预测平均功率 & kW \\
			$\ell_t$ & 第$t$时段小区负载所需电量，$\ell_t=L_t\Delta t$ & kWh \\
			$v_t$ & 第$t$时段光伏预测可用电量，$v_t=\widehat{P}_t\Delta t$ & kWh \\
			$g_t$ & 第$t$时段从外网购入的计划电量 & kWh \\
			$c_t$ & 第$t$时段微网送入电池的外部电量 & kWh \\
			$d_t$ & 第$t$时段电池向微网交付的外部电量 & kWh \\
			$S_t$ & 第$t$时段结束时电池的内部储电量 & kWh \\
			$z_t$ & 充放电状态变量，$z_t=1$表示充电，$z_t=0$表示放电 & 0--1 \\
			$\eta_c$ & 电池充电效率 & -- \\
			$\eta_d$ & 电池放电效率 & -- \\
			$S_{\min}$ & 电池允许的最小储电量 & kWh \\
			$S_{\max}$ & 电池允许的最大储电量 & kWh \\
			$S_0,\ S_T$ & 日初和日末电池储电量 & kWh \\
			$P_c^{\max}$ & 电池最大充电功率 & kW \\
			$P_d^{\max}$ & 电池最大放电功率 & kW \\
			$C_{\max}$ & 单个时段最大外部充电量，$C_{\max}=P_c^{\max}\Delta t$ & kWh \\
			$D_{\max}$ & 单个时段最大外部放电量，$D_{\max}=P_d^{\max}\Delta t$ & kWh \\
			$F$ & 全天总购电费用 & 元 \\
			\bottomrule 
		\end{tabular}
	\end{table}
	\par
	
	\section{数据的整理和预处理}
	
	\section{模型建立与求解}
	\subsection{问题一}
	\subsubsection{建模思路与时间离散化}
	问题一中的电价、小区负载和光伏预测均为已知量，目标是在满足小区负载需求和储能设备
	安全约束的前提下，确定全天各时段的购电量以及储能充放电量，使全天购电费用最小。
	由于储能设备的当前动作会影响后续时段的可用容量，且同一时段不能同时充电和放电，
	该问题同时包含连续变量和0--1变量，因此本文将其建模为混合整数线性规划问题。\par
	为了刻画电价和负载的逐时变化，将一天等分为$T=144$个时段，每个时段长度为
	\begin{equation}
		\Delta t=\frac{24}{144}=\frac{1}{6}\ \mathrm{h}.
	\end{equation}
	第$t$个时段记为$[(t-1)\Delta t,t\Delta t]$。附件中的时间标签作为该时段的结束时刻，
	每条功率数据表示对应时段内的平均功率。设第$t$时段的小区负载平均功率为$L_t$，
	光伏预测平均功率为$\widehat P_t$，则负载电量和光伏预测电量分别为
	\begin{equation}
		\ell_t=L_t\Delta t,\qquad v_t=\widehat P_t\Delta t.
	\end{equation}
	
	\subsubsection{决策变量与目标函数}
	设第$t$时段从外网购入的计划电量为$g_t$，微网送入电池的外部电量为$c_t$，
	电池向微网交付的外部电量为$d_t$，时段结束时电池的内部储电量为$S_t$。
	引入0--1变量$z_t$表示充放电状态：
	\begin{equation}
		z_t=
		\begin{cases}
			1, & \text{第$t$时段充电},\\
			0, & \text{第$t$时段放电或闲置}.
		\end{cases}
	\end{equation}
	由于购电量$g_t$的单位为kWh，目标函数可直接写为各时段购电量与购电价格乘积之和：
	\begin{equation}
		\min F=\sum_{t=1}^{T}p_tg_t.
	\end{equation}
	
	\subsubsection{供需平衡与储能约束}
	在光伏全部消纳且不向外网售电的假设下，每个时段获得的电量等于负载消耗与电池充电量之和，
	因此逐时供需平衡约束为
	\begin{equation}
		g_t+v_t+d_t=\ell_t+c_t,\qquad t=1,2,\ldots,T.
	\end{equation}
	整理可得
	\begin{equation}
		g_t=\ell_t-v_t+c_t-d_t.
	\end{equation}
	由该式可以看出，电池充电会增加购电需求，电池放电会减少购电需求，光伏发电则会直接抵扣
	部分负载电量。\par
	采用电池内部电量计算储电量，充放电效率分别为$\eta_c$和$\eta_d$，所得储电量递推关系为
	\begin{equation}
		S_t=S_{t-1}+\eta_cc_t-\frac{d_t}{\eta_d},
		\qquad t=1,2,\ldots,T.
	\end{equation}
	其中，$\eta_c=\eta_d=0.9$。电池运行过程中必须满足储电量上下限：
	\begin{equation}
		S_{\min}\leq S_t\leq S_{\max},
		\qquad S_{\min}=1200,\quad S_{\max}=10800.
	\end{equation}
	同时，题目要求日初和日末储电量相同，因此
	\begin{equation}
		S_0=S_T=6000.
	\end{equation}
	设电池的最大充电功率和最大放电功率分别为$P_c^{\max}=P_d^{\max}=5000$ kW，
	则每个时段的最大外部充放电量为
	\begin{equation}
		C_{\max}=P_c^{\max}\Delta t=\frac{2500}{3}\ \mathrm{kWh},\qquad
		D_{\max}=P_d^{\max}\Delta t=\frac{2500}{3}\ \mathrm{kWh}.
	\end{equation}
	利用二元变量$z_t$建立充放电功率约束和互斥约束：
	\begin{equation}
		0\leq c_t\leq C_{\max}z_t,
	\end{equation}
	\begin{equation}
		0\leq d_t\leq D_{\max}(1-z_t),
	\end{equation}
	\begin{equation}
		g_t\geq0,\qquad z_t\in\{0,1\}.
	\end{equation}
	当$z_t=1$时，$d_t=0$，电池只能充电；当$z_t=0$时，$c_t=0$，电池只能放电或闲置，
	从而避免同一时段出现同时充放电的现象。
	
	\subsubsection{混合整数线性规划模型}
	综合以上目标函数和约束条件，问题一的完整数学模型为
	\begin{equation}
		\begin{aligned}
			\min\quad & F=\sum_{t=1}^{T}p_tg_t\\
			\text{s.t.}\quad
			& g_t+v_t+d_t=\ell_t+c_t, && t=1,\ldots,T,\\
			& S_t=S_{t-1}+\eta_cc_t-\frac{d_t}{\eta_d}, && t=1,\ldots,T,\\
			& S_{\min}\leq S_t\leq S_{\max}, && t=0,\ldots,T,\\
			& 0\leq c_t\leq C_{\max}z_t, && t=1,\ldots,T,\\
			& 0\leq d_t\leq D_{\max}(1-z_t), && t=1,\ldots,T,\\
			& g_t\geq0,\quad z_t\in\{0,1\}, && t=1,\ldots,T,\\
			& S_0=S_T=6000.
		\end{aligned}
	\end{equation}
	模型中的目标函数和所有约束均为线性形式，其中$z_t$为整数变量，其余变量为连续变量，
	因此该模型属于混合整数线性规划模型。
	
	\subsubsection{基于HiGHS的模型求解}
	在144个时段下，模型共有721个变量，其中购电量、充电量、放电量和充放电状态变量各144个，
	储电量变量145个；二元变量144个，连续变量577个。将附件1中的电价、小区负载和光伏预测功率
	代入模型后，使用SciPy提供的\texttt{optimize.milp}接口调用HiGHS求解器进行求解。\par
	求解器在约0.033秒内收敛，最优性相对间隙为0，求解状态为最优。为排除“求解器报告成功但解
	不满足物理约束”的情况，本文对所得解进行了独立校验，包括逐时供需平衡、储电量递推、
	储电量边界、功率上限、充放电互斥和日循环关系。最大能量残差约为$5.46\times10^{-12}$ kWh，
	所有约束均在允许误差范围内成立。
	
	\subsubsection{最优购电策略与结果分析}
	问题一的最终调度结果如下。全天购电量为59482.698998 kWh，全天购电费为35126.948589元。
	日初和日末储电量均为6000 kWh，全天充电量为20740.666132 kWh，放电量为16799.939567 kWh，
	储能损耗为3940.726565 kWh。\par
	题目指定10分钟时段的购电量如表~\ref{tab:q1_tenminute}所示。
	\begin{table}[htbp]
		\centering
		\caption{问题一指定时段购电量}
		\label{tab:q1_tenminute}
		\begin{tabular}{cc}
			\toprule
			时段 & 购电量（kWh）\\
			\midrule
			10:00--10:10 & 0.000000\\
			12:00--12:10 & 480.412450\\
			14:00--14:10 & 0.000000\\
			16:00--16:10 & 445.431650\\
			18:00--18:10 & 531.894050\\
			20:00--20:10 & 0.000000\\
			\midrule
			全天购电量 & 59482.698998\\
			全天购电费（元） & 35126.948589\\
			\bottomrule
		\end{tabular}
	\end{table}
	
	储能设备六个四小时区间的充放电结果如表~\ref{tab:q1_battery}所示。
	\begin{table}[htbp]
		\centering
		\caption{问题一四小时区间充放电量与起止储电量}
		\label{tab:q1_battery}
		\begin{tabular}{ccc}
			\toprule
			时段 & 充电量（kWh） & 放电量（kWh）\\
			\midrule
			00:00--04:00 & 4500.000000 & 0.000000\\
			04:00--08:00 & 833.333333 & 6365.841200\\
			08:00--12:00 & 4787.964288 & 1702.996983\\
			12:00--16:00 & 5286.035177 & 91.101383\\
			16:00--20:00 & 0.000000 & 5780.131883\\
			20:00--24:00 & 5333.333333 & 2859.868117\\
			\midrule
			0:00储电量（kWh） & \multicolumn{2}{c}{6000.000000}\\
			24:00储电量（kWh） & \multicolumn{2}{c}{6000.000000}\\
			\bottomrule
		\end{tabular}
	\end{table}
	
	从调度结果可以看出，储能设备在夜间低电价及光伏富余时段主要以充电为主，
	在电价较高或负载较大的时段释放电能，从而减少高价时段的外网购电量。
	全天满足
	\begin{equation}
		\sum_{t=1}^{T}d_t=\eta_c\eta_d\sum_{t=1}^{T}c_t=0.81\sum_{t=1}^{T}c_t,
	\end{equation}
	说明充放电过程产生的电能损耗已经计入模型。同时，全天购电量满足
	\begin{equation}
		\sum_{t=1}^{T}g_t
		=\sum_{t=1}^{T}\ell_t-\sum_{t=1}^{T}v_t
		+(1-\eta_c\eta_d)\sum_{t=1}^{T}c_t,
	\end{equation}
	即全天购电量由“负载电量减光伏电量”与储能损耗共同组成。以上结果验证了模型能够在满足
	全部物理约束的前提下给出经济性较好的储能调度方案，完整的144时段计划已保存至
	\texttt{result1.xlsx}。
	\subsection{问题二}
	\subsubsection{建模思路与总体框架}
	问题二每天的电价规律相同，但负载和光伏实际功率随时间变化。0:00制定计划时只能利用此前
	已经掌握的历史信息，因此本文将问题二分解为“滚动预测—风险修正—日前优化—实际运行”四个阶段。
	设第$k$天包含$T=144$个时段，负载和光伏预测功率分别为$\widehat L_{k,t}$和
	$\widehat{PV}_{k,t}$，对应的预测净负荷电量为
	\begin{equation}
		\widehat N_{k,t}
		=\left(\widehat L_{k,t}-\widehat{PV}_{k,t}\right)\Delta t.
	\end{equation}
	为避免仅按预测均值安排购电造成高额紧急购电费用，本文利用历史预测误差的80\%分位数对
	预测净负荷进行修正。修正后的计划净负荷作为当日储能调度MILP的输入，求得的计划购电曲线
	在0:00锁定。当天运行阶段再根据实际负载和光伏执行储能动作，并对电池放电后仍存在的缺口
	进行紧急购电。
	
	\subsubsection{数据整理与滚动预测设置}
	附件2给出2025年每天144个10分钟粒度的实际负载和实际光伏功率。本文将每天的数据按同一
	时段对齐，形成长度为365的日内时间序列。为满足问题二“0:00制定计划”的信息要求，所有预测
	均采用逐日滚动方式：预测第$k$天时只使用第$k$天以前的数据；当天结束后，再将该天实际值
	加入历史序列并预测下一天。\par
	2025年1月用于模型参数选择、历史预测残差估计以及储能设备从初始状态到正式计算期的热启动。
	正式预测和费用统计区间为2025年2月1日至12月31日，共334天。模型输出逐日负载和光伏预测值，
	并整理为\texttt{pred.xlsx}，作为第二问购电与储能调度的输入。
	
	\subsubsection{负载与光伏时间序列预测模型}
	住宅小区负载具有明显的日内周期和周周期。对目标日第$t$个时段，负载模型使用前1、2、3、7日
	的滞后值、前3日和7日的同时刻均值、最近两个相同星期几的均值，以及星期周期、年内周期和
	日内周期特征。设目标时段的特征向量为$X^L_{k,t}$，则负载采用最近31天窗口的岭回归预测：
	\begin{equation}
		\widehat L_{k,t}=X^L_{k,t}\beta_L,
	\end{equation}
	\begin{equation}
		\widehat\beta_L
		=\arg\min_{\beta}
		\left\|
		Y_L-X^L\beta
		\right\|_2^2
		+\lambda_L\|\beta\|_2^2.
	\end{equation}
	其中，岭参数$\lambda_L$在$\{0.1,1,10,100\}$中利用1月数据选择，以平均绝对误差最小为目标，
	当前数据下取$\lambda_L=100$。\par
	光伏发电除受近期天气影响外，还具有明显的太阳高度角季节规律。定义第$k$天的太阳赤纬为
	\begin{equation}
		\delta_k
		=23.44^\circ
		\sin\left[
		\frac{2\pi(D_k-81)}{365}
		\right],
	\end{equation}
	其中$D_k$为年内序号。光伏特征包括前1、2、3、7、14日的滞后值，前3日、7日和14日的同时刻
	均值，以及$\sin(\delta_k)$、$\cos(\delta_k)$和日内周期项。采用扩展窗口岭回归得到
	$\widehat{PV}^{\,ridge}_{k,t}$，并与最近7日同时刻均值
	\begin{equation}
		\widehat{PV}^{\,7}_{k,t}
		=\frac{1}{7}\sum_{j=1}^{7}PV_{k-j,t}
	\end{equation}
	进行融合：
	\begin{equation}
		\widehat{PV}_{k,t}
		=\max\left\{
		0,\ 
		\alpha\widehat{PV}^{\,7}_{k,t}
		+(1-\alpha)\widehat{PV}^{\,ridge}_{k,t}
		\right\}.
	\end{equation}
	利用1月白天时段数据选择参数，当前取岭参数1、融合权重$\alpha=0.4$。将滚动预测结果与
	附件2实际值进行比较，负载和光伏的$R^2$分别为0.9745和0.9899，说明模型能够较好刻画
	全年负载和光伏的变化趋势。
	
	\subsubsection{预测误差与80\%分位数风险修正}
	实际净负荷电量为
	\begin{equation}
		N_{k,t}
		=\left(L_{k,t}-PV_{k,t}\right)\Delta t,
	\end{equation}
	预测误差为
	\begin{equation}
		\varepsilon_{k,t}
		=N_{k,t}-\widehat N_{k,t}.
	\end{equation}
	当$\varepsilon_{k,t}>0$时，表示实际净负荷高于预测，容易产生紧急购电；当
	$\varepsilon_{k,t}<0$时，表示可能出现计划购电富余。由于紧急购电价格为正常电价的5倍，
	本文使用最近60天同一时段的历史预测误差，取80\%分位数作为风险裕量：
	\begin{equation}
		q_{k,t}
		=Q_{0.8}
		\left\{
		\varepsilon_{k-1,t},
		\varepsilon_{k-2,t},
		\ldots,
		\varepsilon_{k-60,t}
		\right\}.
	\end{equation}
	最终用于制定计划的净负荷为
	\begin{equation}
		N^{\mathrm{plan}}_{k,t}
		=\widehat N_{k,t}+q_{k,t}.
	\end{equation}
	该修正会在预测误差偏正时提高计划购电量，在误差偏负时适当降低备购量，从而在计划购电成本
	和五倍紧急购电风险之间取得平衡。
	
	\subsubsection{日前购电与储能调度MILP模型}
	令$g_{k,t}$、$c_{k,t}$、$d_{k,t}$、$u_{k,t}$和$S_{k,t}$分别表示第$k$天第$t$个
	时段的计划购电量、充电量、放电量、计划富余量和时段末储电量，$z_{k,t}$为充放电状态二元变量。
	每天0:00求解如下混合整数线性规划：
	\begin{equation}
		\begin{aligned}
			\min\quad
			& \sum_{t=1}^{T}p_tg_{k,t}\\
			\text{s.t.}\quad
			& g_{k,t}+d_{k,t}-c_{k,t}-u_{k,t}
			=N^{\mathrm{plan}}_{k,t},
			&& t=1,\ldots,T,\\
			& S_{k,t}
			=S_{k,t-1}+\eta_cc_{k,t}-\frac{d_{k,t}}{\eta_d},
			&& t=1,\ldots,T,\\
			& S_{\min}\leq S_{k,t}\leq S_{\max},
			&& t=1,\ldots,T,\\
			& 0\leq c_{k,t}\leq C_{\max}z_{k,t},
			&& t=1,\ldots,T,\\
			& 0\leq d_{k,t}\leq D_{\max}(1-z_{k,t}),
			&& t=1,\ldots,T,\\
			& g_{k,t}\geq0,\qquad u_{k,t}\geq0,
			&& t=1,\ldots,T,\\
			& z_{k,t}\in\{0,1\},
			&& t=1,\ldots,T,\\
			& S_{k,0}=S_{k,T}=S_k^{\mathrm{start}}.
		\end{aligned}
	\end{equation}
	其中$\eta_c=\eta_d=0.9$，$S_{\min}=1200$ kWh，$S_{\max}=10800$ kWh，
	$C_{\max}=D_{\max}=5000\Delta t=833.33$ kWh，
	$S_k^{\mathrm{start}}$为当天计划开始时电池的实际储电量。
	
	\subsubsection{实际运行与跨日滚动求解}
	计划购电曲线在0:00锁定后，当天按实际负载和光伏逐时段运行。定义电池动作前的供需差额为
	\begin{equation}
		a_{k,t}
		=g_{k,t}+PV_{k,t}\Delta t-L_{k,t}\Delta t.
	\end{equation}
	当$a_{k,t}\geq0$时，计划电量和光伏已经满足负载，此时不进行紧急购电，富余电量优先充电：
	\begin{equation}
		c_{k,t}
		=\min\left\{
		a_{k,t},\ 
		C_{\max},\ 
		\frac{S_{\max}-S_{k,t-1}^{\mathrm{act}}}{\eta_c}
		\right\},
	\end{equation}
	未利用富余电量为$u_{k,t}=a_{k,t}-c_{k,t}$。当$a_{k,t}<0$时出现供电缺口，电池优先放电：
	\begin{equation}
		d_{k,t}
		=\min\left\{
		-a_{k,t},\ 
		D_{\max},\ 
		\eta_d\left(S_{k,t-1}^{\mathrm{act}}-S_{\min}\right)
		\right\},
	\end{equation}
	若放电后仍有缺口，则进行紧急购电：
	\begin{equation}
		e_{k,t}
		=\max\left\{0,\ -a_{k,t}-d_{k,t}\right\}.
	\end{equation}
	实际储电量按
	\begin{equation}
		S_{k,t}^{\mathrm{act}}
		=S_{k,t-1}^{\mathrm{act}}
		+\eta_cc_{k,t}
		-\frac{d_{k,t}}{\eta_d}
	\end{equation}
	更新，并满足跨日连续约束
	\begin{equation}
		S_{k+1,0}^{\mathrm{act}}=S_{k,T}^{\mathrm{act}}.
	\end{equation}
	2025年1月1日初始储电量为6000 kWh，经过1月热启动后，2月1日初始储电量达到10800 kWh。
	
	\subsubsection{模型求解结果与分析}
	使用前述滚动预测、80\%分位数风险修正和逐日MILP模型，对2025年2月1日至12月31日进行
	求解。预测精度和全年费用汇总分别如表~\ref{tab:q2_prediction}和表~\ref{tab:q2_total}所示。
	\begin{table}[htbp]
		\centering
		\caption{第二问负载与光伏预测精度}
		\label{tab:q2_prediction}
		\begin{tabular}{lccc}
			\toprule
			预测对象 & $R^2$ & MAE（kW） & RMSE（kW）\\
			\midrule
			小区负载 & 0.9745 & 166.4431 & 224.6384\\
			光伏发电 & 0.9899 & 154.7734 & 301.7602\\
			光伏白天时段 & 0.9735 & 300.4588 & 426.2253\\
			\bottomrule
		\end{tabular}
	\end{table}
	\begin{table}[htbp]
		\centering
		\caption{第二问2025年全年购电与储能结果}
		\label{tab:q2_total}
		\begin{tabular}{lr}
			\toprule
			指标 & 结果\\
			\midrule
			计划购电量 & 22046434.834858 kWh\\
			计划购电费 & 13454464.390255 元\\
			实际充电量 & 6734511.639971 kWh\\
			实际放电量 & 5455069.040090 kWh\\
			紧急购电量 & 126347.963398 kWh\\
			紧急购电费 & 769998.124403 元\\
			总费用 & 14224462.514658 元\\
			2025-12-31期末储电量 & 10672.653652 kWh\\
			\bottomrule
		\end{tabular}
	\end{table}
	结果显示，负载和光伏模型的决定系数均超过0.97，预测值能够较好地跟踪实际变化。
	80\%分位数风险裕量使计划购电量略高于预测均值，但全年紧急购电量仅为126347.963398 kWh，
	只占计划购电量的约0.57\%，有效抑制了五倍电价下的缺电成本。最终总购电费用为
	14224462.514658元，完整的逐日计划购电量、四小时充放电量和紧急购电时段已保存至
	\texttt{result2.xlsx}。
	
	\subsection{问题三}
	\subsubsection{多阶段决策与信息结构}
	第三问在一天中设置四个预报发布和决策时刻：
	\begin{equation}
		\mathcal E=\{0,6,12,18\}.
	\end{equation}
	0:00根据当时的预测制定全天计划购电量，6:00、12:00和18:00根据新预报、已经发生的实际数据
	以及当前储电量，调整尚未结束时段的购电量。最终执行的调整购电量$g_{d,t}^{\mathrm{adj}}$
	取该时段开始前最后一次有效调整值。若时刻$\tau$对应的时段编号为
	\begin{equation}
		k_\tau=6\tau,
	\end{equation}
	则6:00只能调整第37--144个时段，12:00只能调整第73--144个时段，18:00只能调整
	第109--144个时段。时刻$\tau$的信息集记为$\mathcal F_{d,\tau}$，其中包含历史数据、
	当天已经观测到的实际值、附件3在当前时刻发布的光伏预报、当前储电量和电价。任一时刻的决策
	均不能使用尚未发布或尚未发生的信息。
	\iffalse
	对于先验概率，问题二仅考虑了孕妇BMI一项指标，本次将先验概率进行进一步优化，综合考虑把孕妇的各项指标，得到基于多因素的二倍体，三倍体的先验概率 $$P_{MT} , P_{MD}$$ 
	\noindent \textbf{后验估计优化}\par
	为了综合考虑孕妇抽血化验时的各项指标，GC值，z值，读段数等指标，改进原先的极大似然值估计函数，得到基于多因素的极大似然值估计函数 $$Z_{MT} , Z_{MD}$$ 
	
	\fi
	
	\subsubsection{附件3预报的10分钟展开}
	附件3提供整点光伏预报，而购电和储能决策的时间粒度为10分钟。记附件3在时刻$\tau$发布的
	未来第$h$小时光伏预测功率为$\widehat{PV}_{d,\tau,h}$，其对应的10分钟时段编号为
	\begin{equation}
		s_h=6(\tau+h-1),\qquad h=1,2,\ldots,24.
	\end{equation}
	对于位于相邻两个整点节点之间的时段$t$，采用线性插值：
	\begin{equation}
		\widehat{PV}_{d,\tau,t}
		=\widehat{PV}_{d,\tau,h}
		+\frac{t-s_h}{6}
		\left[
		\widehat{PV}_{d,\tau,h+1}
		-\widehat{PV}_{d,\tau,h}
		\right].
	\end{equation}
	曲线两端采用水平外推，即$t\leq s_1$和$t\geq s_{24}$时分别取首尾整点预报值。
	由此得到每个决策时刻可直接用于优化模型的光伏功率曲线。
	
	\subsubsection{0:00计划购电模型}
	附件3不提供负载预报，因此第三问负载仍采用第二问的滚动时间序列模型预测：
	\begin{equation}
		\widehat L_{d,t|0}.
	\end{equation}
	光伏采用附件3在0:00发布的光伏预报，经插值后记为$\widehat{PV}_{d,0,t}$。
	0:00的预测净负荷电量为
	\begin{equation}
		\widehat N_{d,t|0}
		=\left(
		\widehat L_{d,t|0}
		-\widehat{PV}_{d,0,t}
		\right)\Delta t.
	\end{equation}
	令$g_{d,t}^{\mathrm{plan}}$、$c_{d,t}^{\mathrm{plan}}$、
	$d_{d,t}^{\mathrm{plan}}$、$u_{d,t}^{\mathrm{plan}}$、
	$S_{d,t}^{\mathrm{plan}}$和$z_{d,t}$分别表示计划购电量、计划充电量、计划放电量、
	计划富余量、计划储电量和充放电状态。0:00计划模型为
	\begin{equation}
		\begin{aligned}
			\min\quad
			& \sum_{t=1}^{T}p_tg_{d,t}^{\mathrm{plan}}\\
			\text{s.t.}\quad
			& g_{d,t}^{\mathrm{plan}}+d_{d,t}^{\mathrm{plan}}
			-c_{d,t}^{\mathrm{plan}}-u_{d,t}^{\mathrm{plan}}
			=\widehat N_{d,t|0},\\
			& S_{d,t}^{\mathrm{plan}}
			=S_{d,t-1}^{\mathrm{plan}}
			+0.9c_{d,t}^{\mathrm{plan}}
			-\frac{d_{d,t}^{\mathrm{plan}}}{0.9},\\
			& 1200\leq S_{d,t}^{\mathrm{plan}}\leq10800,\\
			& 0\leq c_{d,t}^{\mathrm{plan}}
			\leq\frac{2500}{3}z_{d,t},\\
			& 0\leq d_{d,t}^{\mathrm{plan}}
			\leq\frac{2500}{3}(1-z_{d,t}),\\
			& g_{d,t}^{\mathrm{plan}}\geq0,\qquad
			u_{d,t}^{\mathrm{plan}}\geq0,\qquad
			z_{d,t}\in\{0,1\}.
		\end{aligned}
	\end{equation}
	其中初始储电量为当天0:00的实际储电量：
	\begin{equation}
		S_{d,0}^{\mathrm{plan}}=S_{d,0}^{\mathrm{actual}}.
	\end{equation}
	单日计划不强制日末储电量等于日初值，跨日价值通过后续滚动优化和实际储能连续传递体现。
	
	\subsubsection{6:00、12:00和18:00的调整模型}
	对于任意调整时刻$\tau\in\{6,12,18\}$，只对尚未发生的时段
	$\mathcal T_\tau=\{k_\tau+1,\ldots,T\}$重新优化。更新后的预测净负荷为
	\begin{equation}
		\widehat N_{d,t|\tau}
		=\left(
		\widehat L_{d,t|\tau}
		-\widehat{PV}_{d,\tau,t}
		\right)\Delta t,
		\qquad t\in\mathcal T_\tau.
	\end{equation}
	未来时段满足
	\begin{equation}
		g_{d,t|\tau}^{\mathrm{adj}}+d_{d,t|\tau}
		-c_{d,t|\tau}-u_{d,t|\tau}
		=\widehat N_{d,t|\tau},
	\end{equation}
	\begin{equation}
		S_{d,t|\tau}=S_{d,t-1|\tau}
		+0.9c_{d,t|\tau}
		-\frac{d_{d,t|\tau}}{0.9},
	\end{equation}
	其余储电量、功率和互斥约束与0:00模型相同。过去时段保持原值：
	\begin{equation}
		g_{d,t|\tau}^{\mathrm{adj}}
		=g_{d,t|\tau^-}^{\mathrm{adj}},
		\qquad t\leq k_\tau.
	\end{equation}
	每个时段最终调整量相对0:00计划量的差额为
	\begin{equation}
		\Delta_{d,t}
		=g_{d,t}^{\mathrm{adj}}-g_{d,t}^{\mathrm{plan}}.
	\end{equation}
	将其分解为调高量和调低量：
	\begin{equation}
		\Delta_{d,t}^{+}
		=\max\left\{g_{d,t}^{\mathrm{adj}}-g_{d,t}^{\mathrm{plan}},0\right\},
	\end{equation}
	\begin{equation}
		\Delta_{d,t}^{-}
		=\max\left\{g_{d,t}^{\mathrm{plan}}-g_{d,t}^{\mathrm{adj}},0\right\}.
	\end{equation}
	调高部分按$1.5p_t$收费，调低部分按$0.5p_t$收取违约费。若同一未来时段经历多次调整，
	均以最终调整量相对原始计划量结算，避免重复计费。
	
	\subsubsection{实际运行与紧急购电}
	最终执行时使用附件2的实际负载和实际光伏。定义电池动作前的供需差额为
	\begin{equation}
		a_{d,t}
		=g_{d,t}^{\mathrm{adj}}
		+PV_{d,t}\Delta t
		-L_{d,t}\Delta t.
	\end{equation}
	当$a_{d,t}\geq0$时，不需要放电和紧急购电，富余电量优先充电：
	\begin{equation}
		c_{d,t}
		=\min\left\{
		a_{d,t},\ 
		\frac{2500}{3},\ 
		\frac{10800-S_{d,t-1}^{\mathrm{act}}}{0.9}
		\right\},
	\end{equation}
	未利用富余电量为$u_{d,t}=a_{d,t}-c_{d,t}$。当$a_{d,t}<0$时，电池优先放电：
	\begin{equation}
		d_{d,t}
		=\min\left\{
		-a_{d,t},\ 
		\frac{2500}{3},\ 
		0.9\left(S_{d,t-1}^{\mathrm{act}}-1200\right)
		\right\},
	\end{equation}
	若电池放电后仍有缺口，则紧急购电量为
	\begin{equation}
		e_{d,t}
		=\max\left\{0,\ -a_{d,t}-d_{d,t}\right\}.
	\end{equation}
	实际储电量按
	\begin{equation}
		S_{d,t}^{\mathrm{act}}
		=S_{d,t-1}^{\mathrm{act}}
		+0.9c_{d,t}
		-\frac{d_{d,t}}{0.9}
	\end{equation}
	更新，并满足跨日连续约束
	\begin{equation}
		S_{d+1,0}^{\mathrm{act}}=S_{d,T}^{\mathrm{act}}.
	\end{equation}
	
	\subsubsection{费用模型与滚动求解}
	计划购电费用、调整费用和紧急购电费用分别为
	\begin{equation}
		C_d^{\mathrm{plan}}
		=\sum_{t=1}^{T}p_tg_{d,t}^{\mathrm{plan}},
	\end{equation}
	\begin{equation}
		C_d^{\mathrm{adj}}
		=\sum_{t=1}^{T}
		\left[
		1.5p_t\Delta_{d,t}^{+}
		+0.5p_t\Delta_{d,t}^{-}
		\right],
	\end{equation}
	\begin{equation}
		C_d^{\mathrm{em}}
		=\sum_{t=1}^{T}5p_te_{d,t}.
	\end{equation}
	当天总费用为
	\begin{equation}
		C_d=C_d^{\mathrm{plan}}+C_d^{\mathrm{adj}}+C_d^{\mathrm{em}}.
	\end{equation}
	0:00阶段以计划费用和未来调整、紧急购电费用的条件期望最小为目标；在6:00、12:00和18:00
	重新求解剩余时段的调整问题。每天按“0:00计划—6:00调整—12:00调整—18:00调整—实际运行”的
	顺序滚动执行，并将当天末储电量传递给下一天。
	
	\subsubsection{模型求解结果与分析}
	按照上述模型对2025年2月1日至12月31日进行滚动求解，全年费用和购电量结果如表\ref{tab:q3_total}所示。
	\begin{table}[htbp]
		\centering
		\caption{第三问2025年全年购电与储能结果}
		\label{tab:q3_total}
		\begin{tabular}{lr}
			\toprule
			指标 & 结果\\
			\midrule
			计划购电量 & 22185050.106831 kWh\\
			调整后购电量 & 22690697.195787 kWh\\
			计划购电费 & 13656380.924449 元\\
			调整费用 & 568814.627038 元\\
			紧急购电量 & 12508.001601 kWh\\
			紧急购电费 & 71327.692171 元\\
			总费用 & 14296523.243658 元\\
			2025-12-31期末储电量 & 5913.447914 kWh\\
			\bottomrule
		\end{tabular}
	\end{table}
	与问题二相比，第三问通过6:00、12:00和18:00的光伏预报更新，将全年紧急购电量由126347.963398 kWh降低到12508.001601 kWh，降幅约为90.1\%，说明多阶段调整能够显著降低高价紧急购电风险。当前参数下调整后购电量高于原计划量，计划费用和调整费用之和略有增加，总费用为14296523.243658元，比问题二增加约72060.729元。因此，增加预报更新时刻能够提高计划对实际光伏变化的适应能力，但其经济收益还取决于调整价格参数和预报精度。完整的计划购电量、调整购电量、充放电量及紧急购电结果已保存至\texttt{result3.xlsx}。
	
	\subsection{问题四}
	
	\section{结果分析}

	\clearpage

	\begin{thebibliography}{2}
	\end{thebibliography}

	\clearpage

	\appendix
	\section*{附录}
	
\end{document}
